\question[4]
Eine Datenbank einer Schule enthält die folgenden Tabellen:



\begin{center}
\begin{minipage}[t]{0.45\textwidth}


Schueler \\

\begin{tabular}{|c|l|}
\hline
SchuelerID & Name \\
\hline
1 & Anna Meier \\
2 & Max Schulz \\
3 & Lea Weber \\
\hline
\end{tabular}
\end{minipage}
\hfill
\begin{minipage}[t]{0.45\textwidth}


Kurs \\

\begin{tabular}{|c|l|c|}
\hline
KursID & Kursname & SchuelerID \\
\hline
10 & Informatik & 1 \\
11 & Mathematik & 1 \\
12 & Physik & 2 \\
\hline
\end{tabular}
\end{minipage}
\end{center}

\begin{enumerate}[label=\alph*)]
\item Bestimme die Primärschlüssel beider Tabellen.
\item Bestimme den Fremdschlüssel.
\item Erläutere die Funktion eines Fremdschlüssels.
\end{enumerate}

\begin{solutionbox}{5cm}
\begin{enumerate}
\item Primärschlüssel:
    \begin{itemize}
    \item Tabelle Schueler: SchuelerID
    \item Tabelle Kurs: KursID
    \end{itemize}

\item Fremdschlüssel:     
    SchuelerID in der Tabelle Kurs

\item Ein Fremdschlüssel verweist auf den Primärschlüssel einer anderen Tabelle.    
    Dadurch können Datensätze verschiedener Tabellen miteinander verknüpft werden.    
    In diesem Beispiel zeigt die SchuelerID in der Tabelle Kurs auf einen Schüler in der Tabelle Schueler.
\end{enumerate}
\end{solutionbox}